\section{Literature Survey}
\label{sec:lit}

This section reviews the established work the study rests on; Table~\ref{tab:lit} condenses
the studies that shaped its design.

\subsection{Motor-imagery BCIs}
Wolpaw et al.\ described a BCI as a channel for communication and control that bypasses
nerves and muscles~\cite{wolpaw2002}; its first users were people with almost complete
paralysis~\cite{kubler2001}. The physiological basis of the motor-imagery variant was laid out
by Pfurtscheller and Lopes da Silva: when a movement is made or merely imagined, the power
of the sensorimotor mu (roughly 8--13\,Hz) and beta (roughly 13--30\,Hz) rhythms drops over
the cortical area that controls the limb, and it often rises above its resting level when
the movement ends~\cite{pfurtscheller1999}. Because the motor cortex is organised by body
part and each hemisphere serves the opposite side of the body, these changes differ between,
say, the left and the right hand, and they can be classified trial by
trial~\cite{pfurtscheller2001,pfurtscheller2006}. A sizeable minority of users never reaches dependable control, a problem known as BCI
inefficiency~\cite{vidaurre2010,ahn2015,lee2019openbmi}, so group means can hide people for
whom a decoder fails; we therefore show results patient by patient.

\subsection{Spatial filters and linear classifiers}
Common spatial patterns (CSP) is the classical feature extractor for motor imagery. It
learns a small set of weighted electrode sums whose output variance is large for one class
and small for the other, and the logarithm of that variance becomes the
feature~\cite{ramoser2000,blankertz2008}. Filter-bank CSP repeats this in several bands~\cite{ang2008fbcsp}. The features usually go
to linear discriminant analysis (LDA) or a support vector machine~\cite{cortes1995}, and with
few trials shrinking the covariance estimates stabilises them~\cite{ledoit2004}. A ten-year review of EEG classifiers concluded that shrinkage LDA and
Riemannian classifiers remain among the safest choices when training data are
scarce~\cite{lotte2018review}. CSP matters for this study for one more reason: it is supervised,
so any test trial that enters the fitting of its filters biases the estimate of accuracy.

\subsection{Riemannian geometry}
A second family summarises each trial by the covariance matrix of its channels. These
matrices are symmetric and positive definite and lie on a curved manifold, so straight-line
distances between them are a poor guide. Barachant et al.\ proposed two classifiers that
respect this geometry: minimum distance to the Riemannian mean (MDM), which assigns a trial
to the class whose mean matrix is closest, and tangent-space classification, which maps the
matrices onto a flat space at a reference point where any linear classifier can be
used~\cite{barachant2012}. Reviews describe them as robust and suited to small
samples~\cite{congedo2017,yger2017}, and they underlie the transfer method tested here.

\subsection{Deep networks}
Convolutional networks can learn the filtering from raw EEG. EEGNet is a compact network
that learns temporal filters, then spatial filters for each of them, and then a separable
summary, with only a few thousand weights; across several BCI paradigms it matched the
reference pipelines~\cite{lawhern2018}. Schirrmeister et al.\ introduced shallow and deep
convolutional networks for EEG decoding~\cite{schirrmeister2017}, the starting point of the
braindecode library used here~\cite{braindecode}. Later designs add attention~\cite{song2023conformer} or pre-train on thousands of hours of
EEG~\cite{kostas2021bendr,jiang2024labram,wang2025cbramod}, but many deep-learning EEG
studies are hard to reproduce~\cite{roy2019}, so we treat EEGNet as one candidate among
several.

\subsection{Calibration and transfer learning}
Lotte surveyed ways to reduce or remove calibration: regularised estimators, adaptation
during use and, above all, reuse of data recorded from other sessions or other
people~\cite{lotte2015}. Early studies pooled many users~\cite{fazli2009}, reused earlier
sessions~\cite{krauledat2008}, adapted the classifier online~\cite{shenoy2006} or learned users
as related tasks~\cite{jayaram2016}. The obstacle is that each person's EEG occupies a different
region of feature space. Riemannian re-centring moves every person's covariance matrices so
that their own mean becomes the identity matrix, after which data from different people can
be pooled~\cite{zanini2018}. Riemannian Procrustes analysis also equalises their dispersion
and rotation~\cite{rodrigues2019}, and a related scheme combined other users' class means
with the new user's in a different BCI paradigm~\cite{kalunga2018}. Euclidean alignment
applies the re-centring idea to the raw trials by whitening them with the inverse square root
of the person's mean covariance, so that a neural network can be trained on the aligned
signals; it needs no labels from the new user~\cite{he2020}. A review found alignment and
fine-tuning to be the most frequent transfer tools in EEG BCIs~\cite{wu2022review}, and
fine-tuned networks have since driven a robotic hand finger by finger~\cite{ding2025}. After stroke, transfer has mainly drawn on the same patient's previous sessions, which
improved accuracy when the current session offered few trials~\cite{giles2022}; borrowing from
previous sessions or other patients also lifted sessions that began poorly~\cite{cao2021}. Nearly all of this work concerns commands that people already
know, such as moving their own hand.

\subsection{BCIs after stroke}
After a stroke one arm is often weak. A meta-analysis of controlled trials found that BCI
therapy improved upper-limb function more than control treatments~\cite{cervera2018}, and in
one trial muscle stimulation triggered by the brain produced gains that outlasted the
treatment~\cite{biasiucci2018}. Other trials paired imagery with a robot~\cite{ang2015rct} or with feedback of the
patient's own hand~\cite{pichiorri2015}.
Stroke also makes decoding harder, because the lesion and therefore the EEG response to
movement differ from one patient to the next~\cite{park2016}, which is exactly what makes
transfer between patients uncertain. Several public stroke datasets have appeared recently,
covering acute stroke~\cite{liu2024acute}, chronic stroke~\cite{lu2026chronic} and imagined
leg movements~\cite{liu2025lowerlimb}; the one used here is, to our knowledge, the only one
with sixth-finger imagery~\cite{wang2026dataset}.

\subsection{Supernumerary limbs and sixth-finger imagery}
Supernumerary robotic limbs add a body part instead of replacing one; reviews cover extra
fingers that help a hand to grasp and extra arms for industrial
tasks~\cite{prattichizzo2021,tong2021}, and Eden et al.\ named the control signal as one of
the central open problems of movement augmentation~\cite{eden2022}. The brain can learn to command an extra digit: people born with six fingers move the extra
one independently~\cite{mehring2019}, a robotic third thumb worn for a few days changed the
hand's representation~\cite{kieliba2021}, and a third robotic arm has been driven by a
brain--machine interface~\cite{penaloza2018}. Wang et al.\ introduced sixth-finger motor imagery. In healthy
people it produced stronger ERD and better classification than imagery of a natural
limb~\cite{wang2024jbhi}; in stroke patients it produced stronger ERD and functional
connectivity~\cite{wang2024tnsre}. Related work described its EEG features~\cite{liu2022jne} and resting-state changes after
training~\cite{liu2022frontiers}. None of
these studies asks how many trials a decoder needs to learn the new command, or whether the
command can still be separated from imagery of the patient's own hand when the evaluation is
strict.

\subsection{Pitfalls in evaluation}
An accuracy figure is meaningful only if the decoder never learned from the trials it is
scored on. Leakage, any path by which test information reaches training, was found across
many fields that apply machine learning and made results irreproducible~\cite{kapoor2023};
every supervised step, including spatial filters, feature selection and scaling, must be
fitted within the training folds~\cite{varoquaux2017,lemm2011}. Splitting short segments of
one recording across training and test lets a decoder recognise the recording instead of
the brain state, and in clinical EEG it inflated accuracy greatly~\cite{brookshire2024}.
When all trials of a class are recorded in one block, decoders can learn slow drifts
instead of the task~\cite{li2021perils}. Chance on a small test set must come from the binomial distribution~\cite{combrisson2015},
and many comparisons call for rank tests with a correction~\cite{demsar2006}. Shared benchmarks such as MOABB fix the evaluation across
datasets~\cite{jayaram2018moabb,chevallier2024moabb}; two of the datasets it serves, the
four-class set of BCI Competition~IV~\cite{tangermann2012} and the set of Cho et
al.~\cite{cho2017}, serve here as healthy controls for the leakage measurement.

\begin{table}[!t]
  \centering
  \caption{Earlier studies that shaped the design, and what each leaves open for the
  question of this paper.}
  \label{tab:lit}
  \footnotesize
  \renewcommand{\arraystretch}{1.15}
  \begin{tabularx}{\textwidth}{@{}>{\raggedright\arraybackslash}p{2.7cm}>{\raggedright\arraybackslash}p{3.3cm}LL@{}}
    \toprule
    Study & Population and method & Finding & Open for this question \\
    \midrule
    Barachant et al.~\cite{barachant2012} & public healthy datasets; MDM and tangent space & as accurate as CSP pipelines with less processing & each user calibrated alone \\
    Zanini et al.~\cite{zanini2018} & healthy BCI data; Riemannian re-centring & pooled data of different people reduce calibration & familiar commands, healthy users \\
    Giles et al.~\cite{giles2022}; Cao et al.~\cite{cao2021} & stroke patients; transfer across sessions & earlier sessions help when current data are scarce or poor & a patient's own hand, not a new command \\
    Wang et al.~\cite{wang2024jbhi,wang2024tnsre} & healthy people, then stroke patients; ERD and classification & sixth-finger imagery gives strong ERD & calibration not studied \\
    Wang et al.~\cite{wang2026dataset} & \nSubjects{} stroke patients, three sessions; CSP with LDA or SVM & \rqOnePublishedGroupBCspLda--\rqOnePublishedGroupTCspLda\,\% for hand against sixth finger before therapy & filters fitted on all trials before the split; one centre \\
    Kapoor and Narayanan~\cite{kapoor2023}; Brookshire et al.~\cite{brookshire2024} & many fields; clinical EEG & leakage and segment-wise splits inflate results & not measured on motor imagery of a new command \\
    \bottomrule
  \end{tabularx}
\end{table}

In summary, the decoders and the calibration shortcuts available today were developed and
validated mostly with healthy volunteers and movements they already know. The one public
sixth-finger dataset reports high accuracy, yet with a fitted step placed before the split
that the evaluation literature warns against. Three questions follow, and this paper takes
them up in order: how well hand and sixth-finger imagery separate under a leakage-free
evaluation; how much calibration a stroke patient needs for the new command, with and
without other patients' data; and whether the sixth finger behaves as a command of its own
or only signals that some movement is being imagined.
